%%% -*- mode: LaTeX -*-
%
% Copyright © 2012-2016 Philipp Büttgenbach
%
% ulphi.en.tex is part of ulphi, a CAE tool and library for computing
% electromagnetic fields.
%
% All rights reserved.
%
% This program is distributed in the hope that it will be useful, but
% WITHOUT ANY WARRANTY; without even the implied warranty of
% MERCHANTABILITY or FITNESS FOR A PARTICULAR PURPOSE.
%

\NeedsTeXFormat{LaTeX2e}

\documentclass[25pt,a0paper,innermargin=18mm]{tikzposter}

\geometry{papersize={847mm,1195mm}}

% -*- coding: utf-8 -*-
%
% Copyright © 2012-2015 Philipp Büttgenbach
%
% This file is part of ulphi, a CAE tool and library for computing
% electromagnetic fields.
%
% Permission is granted to copy, distribute and/or modify this
% document under the terms of the \gls{gnu} Free Documentation
% License, Version 1.3 or any later version published by the Free
% Software Foundation; with no Invariant Sections, no Front-Cover
% Texts, and no Back-Cover Texts.  A copy of the license is included
% in the section entitled GNU Free Documentation License.
%

%\usepackage[utf8]{inputenx}
\usepackage[ngerman,main=english]{babel}
\usepackage{amsmath,amsthm,amssymb}
\usepackage[math-style=ISO,bold-style=ISO,nabla=upright,partial=upright,vargreek-shape=unicode]{unicode-math}
\usepackage[squaren]{SIunits}
\usepackage{graphicx}
\usepackage{pgfplots}
\usepackage[style=authoryear]{biblatex}

\newcommand*{\personname}[1]{\textsc{#1}}
\newcommand*{\partdiff}[2]{\ensuremath{\frac{\partial#1}{\partial#2}}}
\newcommand*{\aninteg}[4]{\ensuremath{#1_{#2}#3\;\mathrm{d}#4}}
\newcommand*{\integ}[3]{\aninteg{\int}{#1}{#2}{#3}}
\newcommand*{\integg}[4]{\ensuremath{\int_{#1}^{#2}#3\;\mathrm{d}#4}}
\newcommand*{\iinteg}[3]{\aninteg{\iint}{#1}{#2}{#3}}
\newcommand*{\iiinteg}[3]{\aninteg{\iiint}{#1}{#2}{#3}}
\newcommand*{\ointeg}[3]{\aninteg{\oint}{#1}{#2}{#3}}
\newcommand*{\oiinteg}[3]{\aninteg{\oiint}{#1}{#2}{#3}}
\newcommand*{\vect}[1]{\ensuremath{\symbf{#1}}}
\newcommand*{\matr}[1]{\vect{#1}}
\newcommand*{\jacmatr}{\ensuremath{\matr{G}_{\mathrm{J}}}}
\newcommand*{\xmax}[1]{\ensuremath{#1_{\text{max}}}}
\newcommand*{\diag}{\ensuremath{\operatorname{diag}}}
\newcommand*{\vectnorm}[1]{\ensuremath{\lvert#1\rvert}}
\newcommand*{\laplaceop}{\ensuremath{\operatorname{\increment}}}
\newcommand*{\grad}{\ensuremath{\operatorname{grad}}}
\newcommand*{\divop}{\ensuremath{\operatorname{div}}}

\bibliography{../include/literature}

\pgfplotsset{
  compat=1.8,
  grid=major,
  grid style={line width=.18mm},
  tick style={line width=.35mm},
}

\definecolor{copper}{cmyk}{0,.32,.72,.28}

\input{../include/config.tex}

%%% Local Variables: 
%%% mode: LaTeX
%%% ispell-local-dictionary: "en_US"
%%% End:


\usepackage{url}

\usepgfplotslibrary{groupplots}

\newcommand*{\Um}{\ensuremath{U_{\mathrm{m}}}}

 % Title, Author, Institute
\title{\textbf{\sffamily A python module for magnetostatics}}
\author{\sffamily Philipp Büttgenbach}
\institute{\url{https://www.github.com/pbuettgen/ulphi} --- \url{mailto:ulphi@ist-einmalig.de}}
\titlegraphic{\includegraphics[height=11cm]{../images/Ulphi-logo.pdf}}

 %Choose Layout
% \usetheme{Autumn}
% \usetheme{Basic}
 \usetheme{Board}
% \usetheme{Default}
%\usetheme{Desert}
% \usetheme{Envelope}
% \usetheme{Rays}
% \usetheme{Simple}
% \usetheme{Wave}

\begin{document}

\maketitle
 
\begin{columns}

  \column{0.3}

  \block{\sffamily What is it?}{
    \begin{itemize}
    \item Compute 2D magnetic fields
    \item Static and quasistatic fields supported
    \item Based on finite integration technique (FIT)
    \item Targets at models with a rather simple geometry
    \item Couple polar and cartesian sections in one model
    \item Strong support for nonlinear and anisotropic materials
    \end{itemize}
  }

  \block{\sffamily How does it work?}{
    \innerblock{\sffamily Equations in magnetostatics}{
      \begin{align*}
        \label{eq:1}
        \nabla\times\left(\nu\vect{B}\right) &= \vect{J}
        &\Leftrightarrow&&
        \ointeg{l}{\nu\vect{B}}{\vect{l}}&= \Theta \\
        %% 
        \nabla\times\vect{A} &= \vect{B} &\Leftrightarrow&&
        \ointeg{l}{\vect{B}}{\vect{l}}&= \Phi
      \end{align*}
    }
    \vspace{1em}
    \innerblock{\sffamily A basic FIT grid}{
      Discretization is done using the finite integration technique
      (FIT). \par \vspace{2ex}
      \begin{minipage}{\linewidth}
        \centering\input{../images/FIT-Grid-3D-Large.pdf_tex}         
      \end{minipage}
    }
    \vspace{1em}
    \innerblock{\sffamily Coupling grids}{
      Grid patches can be coupled even when having different types.
      \begin{itemize}
      \item Smooth transition from cartesian to polar sections
      \item No interpolation of material properties necessary!
      \end{itemize}
      \vspace{1ex}
      \centering\input{../images/Interface-polar-cartesian-Large.pdf_tex} 
    }
  }

  \column{0.7}

  \block{\sffamily Examples}{
    \begin{minipage}[t]{.52\colwidth}
      \innerblock{\sffamily Coaxial line}{
        An analytical solution is
        known and can be used to verify Ulphi:
        \begin{multline*}
          A(0)-A(r) = \frac{\mu_0J_1}{2} \\ \cdot
          \begin{cases}
            r^2/2 & r \le r_1\text{~,} \\
            r_1^2\left(\frac{1}{2}+\ln\frac{r}{r_1}\right)
            & r_1 < r \le r_2\text{~,} \\[.5ex]
            r_1^2\left(\frac{1}{2}+\ln\frac{r_2}{r_1}
              +\left(1+\frac{r_2^2}{r_3^2-r_2^2}\right)\ln\frac{r}{r_2}
              +\frac{r^2-r_2^2}{2(r_3^2-r_2^2)}\right)
            & r_2 < r \le r_3\text{~.}
          \end{cases}
        \end{multline*}

        \begin{minipage}{.7\linewidth}
          \centering
          \begin{tikzpicture}
            \begin{axis}[
              title=Magnetic vector potential (z-component),
              width=.85\linewidth,
              height=.5253288904374106\linewidth,
              xlabel=${r}/{\milli\metre}\quad\longrightarrow$,
              ylabel=$A\cdot\frac{\metre}{\micro\weber}\quad\longrightarrow$,
              xtick={0,5,...,20},
              ytick={0,20,...,100},
              legend style={at={(0.95,0.95)}, anchor=north east}, 
              line width=.35mm]
              \addplot[mark=diamond,only marks] table[x=x, y=NumA]
              {../datasets/coaxial_line.csv}; 
              \addplot[smooth,mark=none,line width=.7mm,style=dotted]
              table[x=x, y=SymA] {../datasets/coaxial_line.csv}; 
              \legend{numerical, analytical}
            \end{axis}
          \end{tikzpicture}
        \end{minipage}
        \begin{minipage}{.275\linewidth}
          \centering
          \begin{tikzpicture}
            \fill[color=copper] (0,0) circle [radius=3cm];
            \fill[white] (0,0) circle [radius=2.5cm];
            \fill[color=copper] (0,0) circle [radius=.75cm];
            \draw (0,0) node {$I_1$};
            \draw (0, -1.5) node {$I_2=0$};
            \draw[line width=.7mm,->] 
            (-.75, -3.25) node[below] {$I_3=-I_1$} -- (-.5,-2.75);
            \draw[line width=.7mm,<-] (45:.75)
            -- (3.25,3.25)
            -- (3.75,3.25) node[near start, above] {$r_1$};
            \draw[line width=.7mm,<-] (21.037511025421818:2.5)
            -- (3.25,1.25)
            -- (3.75, 1.25) node[near start, above] {$r_2$};
            \draw[line width=.7mm,<-] (-31.607502246248906:3)
            -- (3.25,-2)
            -- (3.75, -2) node[near start, above] {$r_3$};
          \end{tikzpicture}
        \end{minipage}
      }
    \end{minipage}
    \begin{minipage}[t]{.42\colwidth}
      \innerblock{\sffamily Simple iron core}{
        Check of average flux density:
        \begin{center}
          \emph{Ulphi:} \unit{1.66}{\tesla} \hfil \emph{Average iron
            path:} \unit{1.59}{\tesla}
        \end{center}
        \begin{itemize}
        \item Ulphi can be trusted
        \end{itemize}
        \vfil \centering
        \begin{tikzpicture}
          \begin{axis}[
            title={$A\cdot\frac{\metre}{\milli\weber}$},
            width=.8\linewidth,
            axis equal={true},
            xmin={0},
            xmax=50,
            ymin=-5,
            ymax=60,
            xtick={-10,0,...,60},
            ytick={-10,0,...,60},
            xlabel=$x/{\milli\metre}\quad\longrightarrow$,
            ylabel=$y/{\milli\metre}\quad\longrightarrow$,
            grid=major,
            grid style={line width=.18mm},
            tick style={line width=.35mm},
            line width=.35mm]
            %% iron core shape
            \addplot[gray,fill=gray!20,opacity=.33,line width=.7mm] 
            coordinates {
              (15,0) (40,0) (40,54) (0, 54) (0, 29) (15,29)} -- cycle;
            %% coils
            \addplot[copper,fill=copper,opacity=.33,line width=.7mm]
            coordinates {(5,0) (10,0) (10,24) (5,24)} -- cycle;
            \addplot[copper,fill=copper,opacity=.33,line width=.7mm]
            coordinates { (45,0) (50,0) (50,24) (45,24)} -- cycle;
            \addplot[contour gnuplot, line width=.7mm]
            table[x=x, y=y, z=potential] {../datasets/simple_iron_core.csv};
          \end{axis}
        \end{tikzpicture}
      }
    \end{minipage}
    \innerblock{\sffamily Wound iron core}{
      \hfil
      \begin{tabular}[b]{cc}
        \input{../images/Wound-iron-core-Large.pdf_tex} &
        \input{../images/wound-iron-core-model-Large.pdf_tex}
        \tabularnewline[1ex]
        Device & Model \vspace{5ex}
      \end{tabular}
      \hfil
      \begin{tikzpicture}
        \begin{groupplot}[
          group style={group size=1 by 2, vertical sep=12ex},
          width=.3\linewidth,
          height=.18541019662496844\linewidth]
          \nextgroupplot[
          title={magnetomotive force},
          xlabel=$t/{\milli\second}\quad\longrightarrow$,
          ylabel=$\Um/{\kilo\ampere}\quad\longrightarrow$,
          xtick={-5,0,...,30},
          legend style={at={(0.95,0.95)},anchor=north east},
          line width=.35mm]
          \addplot[line width=.7mm] table[x=t,
          y=y1]{../datasets/wound_iron_core_time.csv}; 
          \addplot[smooth, line
          width=.7mm, style=dashed] table[x=t,
          y=y2]{../datasets/wound_iron_core_time.csv}; 
          \legend{waveform, fundamental}
          \nextgroupplot[
          title={magnetomotive force (spectrum)},
          xlabel=harmonic order$\quad\longrightarrow$,
          ylabel=$\Um/{\kilo\ampere}\quad\longrightarrow$,
          xtick={-20,-15,...,20},
          ytick={0,.1,...,.5},
          line width=.35mm]
          \addplot[ycomb, line width=.7mm]
          table{../datasets/wound_iron_core_freq.csv};
        \end{groupplot}
      \end{tikzpicture}
      \hfil
    } 
  }

\block{\sffamily Ideas for future development}{
  \begin{itemize}
  \item Permanent magnets
  \item Force (and torque) evaluation
  \item Implement linear Laplace equation (required for electrostatics
    and thermostatics)
  \item High frequency FIT
  \item Graphical user interface
  \end{itemize}
}


\end{columns}

\begin{columns}
  \column{.5}
  \block{\sffamily References}{
    \nocite{Salon:1995}
    \nocite{Silvester:1996}
    \nocite{Weiland:1977}
    \nocite{vanRienen:2010}
    \printbibliography[title={\relax}]
  }
  \column{.5}
  \block{\sffamily Acknowledgement}{
    Thanks a lot to my parents for sponsoring my trip to EuroScipy
    2016 in Erlangen.
  }
  \block{\relax}{
    \hfil
    \includegraphics[height=4cm]{../images/python-cpp-logo.png} \hfil
    \includegraphics[height=4cm]{../images/boost-library-logo.png} \hfil
    \includegraphics[height=4cm]{../images/Eigen-logo_silly-professor.png} \hfil
    \includegraphics[height=4cm]{../images/imprint-logo_Wir-machen-druck.png} \hfil
  }
\end{columns}

\end{document}

%%% Local Variables: 
%%% mode: latex
%%% TeX-master: nil
%%% ispell-local-dictionary: "en_US"
%%% End: 
